\documentclass[11pt]{article}
\usepackage[top=2cm]{geometry}

\title{Homework 15.2: Jim Blinn's Biography}
\author{Alejandro Cornejo Pérez}
\date{May 2026}
\setlength{\parskip}{\baselineskip}
\begin{document}

\maketitle
Jim Blinn is an American computer scientist, he was born in 1949 and is 77 years old. He is considered as one of the most knowledgeable people in the computer graphics. He's known for working in realistic image rendering, shading techniques, and animation. His work helped change modern 3D graphics used in movies, video games, simulations, and scientific visualization.

He studied engineering and computer science during the birth of digital graphics. He later worked at companies like the Jet Propulsion Laboratory. This is where he created computer generated visualizations for space missions. One of his most known contributions was producing animations for NASA’s Voyager missions. These animations helped the public better understand planetary exploration.

Jim Blinn contributed with the creation of Blinn-Phong shading. This is an improved version of the Phong model because it introduced a more efficient method for calculating specular highlights on surfaces. Another major contribution by him was the development of bump mapping in 1978. This technique allows surfaces to appear more detailed and textured without actually increasing the number of polygons. Instead of changing the geometry it just modifies the lighting calculations to create the illusion of depth and roughness.

Jim Blinn also contributed to mathematical visualization and educational content. He became well known for his series called Jim Blinn’s Corner, here he explained complex graphics concepts in a simple and entertaining way. He helped generations of students and professionals understand topics such as rendering, geometry, lighting, and animation in entertaining ways.

His impact is still seen today in companies like Pixar, NVIDIA, and video game studios like Capcom. He laid the bases for realistic rendering techniques and they continue to evolve in products all over the world. This is why Jim Blinn is considered one of the most important figures in the history of computer graphics.
\end{document}
